\documentclass[12pt]{article}
\usepackage[ukrainian]{babel}
\usepackage{fontspec}
% --- НАЛАШТУВАННЯ ШРИФТІВ ---
\setmainfont{Schoolbook-Regular.otf}[
    Path = ../../,
    BoldFont = Schoolbook-Bold.otf,
    ItalicFont = Schoolbook-Italic.otf,
    BoldItalicFont = Schoolbook-BoldItalic.otf
]
\usepackage[italic]{mathastext}
\usepackage[a4paper,margin=1.8cm,bottom=2cm,top=2cm]{geometry}
\usepackage{amsmath,amssymb}
\usepackage{tikz}
\usepackage{pgfplots}
\pgfplotsset{compat=1.16}
\usetikzlibrary{calc,patterns,angles,quotes}
\usepackage{xcolor,array,fancyhdr,enumitem,multicol}
\usepackage{colortbl}
\usepackage{tcolorbox}
\tcbuselibrary{skins}
\usepackage[hidelinks]{hyperref}
\usepackage[firstpage=false, text=@pvtr2525, scale=0.6, color=gray!12]{draftwatermark}
\definecolor{mainGreen}{RGB}{34, 120, 64}
\definecolor{yearOrange}{RGB}{220, 100, 30}
\definecolor{yearcolor}{RGB}{220, 100, 30}
\definecolor{titleBg}{RGB}{225, 240, 220}
\definecolor{instrBg}{RGB}{255, 240, 220}
\definecolor{instrBorder}{RGB}{220, 150, 80}
\pagestyle{fancy}
\fancyhf{}
\renewcommand{\headrulewidth}{0.4pt}
\renewcommand{\footrulewidth}{0.4pt}
\fancyhead[C]{\small\color{gray!80} Збірник НМТ \textendash{} сесії 23.05--5.06.2026}
\fancyhead[R]{\small\color{mainGreen}\textbf{\href{https://t.me/pvtr2525}{@pvtr2525}}}
\fancyfoot[L]{\small\color{gray!80}\href{https://t.me/pvtr2525}{ТГ: @pvtr2525}}
\fancyfoot[C]{\small\color{gray!80}\thepage}
\fancyfoot[R]{\small\color{gray!80}НМТ 2026}
% --- ТАБЛИЦІ ВІДПОВІДЕЙ ---
\newcommand{\answerTable}[5]{%
\par\vspace{0.15cm}
\noindent
\begin{tabular}{|*{5}{>{\centering\arraybackslash}m{3cm}|}}
\hline
\rule[-0.15cm]{0pt}{0.6cm}\textbf{А} & \rule[-0.15cm]{0pt}{0.6cm}\textbf{Б} & \rule[-0.15cm]{0pt}{0.6cm}\textbf{В} & \rule[-0.15cm]{0pt}{0.6cm}\textbf{Г} & \rule[-0.15cm]{0pt}{0.6cm}\textbf{Д} \\
\hline
\rule[-0.2cm]{0pt}{0.75cm}#1 & \rule[-0.2cm]{0pt}{0.75cm}#2 & \rule[-0.2cm]{0pt}{0.75cm}#3 & \rule[-0.2cm]{0pt}{0.75cm}#4 & \rule[-0.2cm]{0pt}{0.75cm}#5 \\
\hline
\end{tabular}
\par\vspace{0.25cm}
}
\newcommand{\answerTableTall}[5]{%
\par\vspace{0.15cm}
\noindent
\begin{tabular}{|*{5}{>{\centering\arraybackslash}m{3cm}|}}
\hline
\rule[-0.2cm]{0pt}{0.7cm}\textbf{А} & \rule[-0.2cm]{0pt}{0.7cm}\textbf{Б} & \rule[-0.2cm]{0pt}{0.7cm}\textbf{В} & \rule[-0.2cm]{0pt}{0.7cm}\textbf{Г} & \rule[-0.2cm]{0pt}{0.7cm}\textbf{Д} \\
\hline
\rule[-0.45cm]{0pt}{1.2cm}#1 & \rule[-0.45cm]{0pt}{1.2cm}#2 & \rule[-0.45cm]{0pt}{1.2cm}#3 & \rule[-0.45cm]{0pt}{1.2cm}#4 & \rule[-0.45cm]{0pt}{1.2cm}#5 \\
\hline
\end{tabular}
\par\vspace{0.25cm}
}
\newcommand{\instructionBox}[1]{%
\par\vspace{0.3cm}
\begin{tcolorbox}[colback=instrBg, colframe=instrBorder, boxrule=0.6pt, arc=2pt,
    left=8pt, right=8pt, top=4pt, bottom=4pt]
\centering\small\bfseries #1
\end{tcolorbox}
\par\vspace{0.15cm}
}
\newcommand{\sectionTitle}[1]{%
\par\vspace{0.4cm}
\begin{tcolorbox}[colback=titleBg, colframe=titleBg, boxrule=0pt, arc=8pt,
    halign=center, fontupper=\Large\bfseries\color{mainGreen}]
#1
\end{tcolorbox}
\par\vspace{0.3cm}
}
\newcommand{\task}[2]{\noindent\textbf{#1.}\ \ #2\par\vspace{0.2cm}}
% --- НМТ-СТИЛЬ ПОЛЕ ВІДПОВІДІ ---
\newcommand{\nmtAnswerBox}{%
\par\vspace{0.2cm}\noindent
Відповідь:\ \
\begin{tabular}{@{}*{4}{|>{\centering\arraybackslash}p{0.8cm}}|@{\hspace{0.1cm}\textbf{,}\hspace{0.1cm}}*{3}{|>{\centering\arraybackslash}p{0.8cm}}|@{}}
\hline
\rule[-0.2cm]{0pt}{0.7cm} & & & & & & \\
\hline
\end{tabular}
\par\vspace{0.3cm}
}
% --- СІТКА ДЛЯ МАТЧИНГУ ---
\newcommand{\matchingGrid}{%
    \begingroup
    \renewcommand{\arraystretch}{1.15}
    \setlength{\tabcolsep}{4pt}
    \footnotesize
    \begin{tabular}{r|c|c|c|c|c|}
         \multicolumn{1}{c}{} & \multicolumn{1}{c}{\textbf{А}} & \multicolumn{1}{c}{\textbf{Б}} & \multicolumn{1}{c}{\textbf{В}} & \multicolumn{1}{c}{\textbf{Г}} & \multicolumn{1}{c}{\textbf{Д}} \\ \cline{2-6}
         \textbf{1} & \phantom{X} & \phantom{X} & \phantom{X} & \phantom{X} & \phantom{X} \\ \cline{2-6}
         \textbf{2} & & & & & \\ \cline{2-6}
         \textbf{3} & & & & & \\ \cline{2-6}
    \end{tabular}
    \endgroup
}
% --- ДОДАТКОВІ МАКРОСИ ЗБІРНИКА ---
\newcommand{\taskBlock}[1]{%
\par\vspace{0.45cm}
\noindent{\color{mainGreen}\rule{\linewidth}{0.8pt}}\par\vspace{0.12cm}
\noindent{\small\bfseries\color{mainGreen}#1}\par\vspace{0.25cm}
}
\newcounter{zad}
\newcommand{\zadtask}[1]{\stepcounter{zad}\noindent\textbf{\thezad.}\ \ #1\par\vspace{0.2cm}}
\newcommand{\zadnum}{\stepcounter{zad}\textbf{\thezad.}\ \ }
\newcounter{ans}
\newcommand{\ansitem}[1]{\stepcounter{ans}\noindent\textbf{\theans.}\ #1\par\vspace{0.07cm}}
\newcommand{\nmtyear}[1]{\hfill{\small\color{yearOrange}(НМТ #1)}}% --- ДОДАТКОВІ МАКРОСИ (розв'язки, зміст) ---
\tcbuselibrary{breakable}
\newcommand{\solution}[1]{%
\par\vspace{0.12cm}
\begin{tcolorbox}[colback=titleBg!45, colframe=mainGreen!55, boxrule=0.4pt, arc=2pt,
    left=6pt, right=6pt, top=3pt, bottom=3pt, breakable]
\footnotesize\textbf{Короткий розв'язок (оригінал).}\ #1
\end{tcolorbox}
\par\vspace{0.12cm}
}
\setcounter{tocdepth}{2}
\newcommand{\chapterTitle}[1]{%
\clearpage\phantomsection\addcontentsline{toc}{section}{#1}%
\sectionTitle{#1}}
\newcommand{\typeTitle}[1]{%
\phantomsection\addcontentsline{toc}{subsection}{\quad #1}%
\par\vspace{0.3cm}
\noindent{\large\bfseries\color{yearOrange}#1}\par\vspace{0.08cm}
\noindent{\color{yearOrange}\rule{\linewidth}{0.4pt}}\par\vspace{0.2cm}}
\newcommand{\ansTheme}[1]{\par\vspace{0.25cm}\noindent{\bfseries\color{mainGreen}#1}\par\vspace{0.12cm}}
\newcommand{\ansType}[1]{\par\vspace{0.1cm}\noindent{\itshape\color{yearOrange}#1}\par\vspace{0.08cm}}
\newcommand{\solbtn}[1]{\hypertarget{task:#1}{}\par\vspace{0.05cm}\noindent\hyperlink{sol:#1}{\fbox{\footnotesize пояснення $\rightarrow$}}\par\vspace{0.2cm}}
\begin{document}
\thispagestyle{empty}
\vspace*{1.2cm}
\begin{center}
{\Large\bfseries\color{mainGreen} ОРИГІНАЛЬНІ ЗАВДАННЯ НМТ-2026}\\[0.4cm]
{\Huge\bfseries\color{mainGreen} РОЗДІЛ 1. ЧИСЛА І ВИРАЗИ}\\[0.6cm]
{\large Лише оригінали сесій 23.05--5.06.2026, без аналогів}\\[1cm]
{\large\color{mainGreen}\textbf{\href{https://t.me/pvtr2525}{Телеграм: @pvtr2525}}}
\end{center}
\clearpage
\chapterTitle{РОЗДІЛ 1. ЧИСЛА І ВИРАЗИ}
\typeTitle{Числа та обчислення}
\taskBlock{Ділення з остачею (порції) --- сесія 30.05, завдання №5}
\zadtask{Для приготування однієї порції коктейлю потрібно $40$ г морозива. Яку \textit{найбільшу} кількість таких порцій коктейлю можна приготувати з $0{,}5$ кг морозива?}
\answerTable{$10$}{$11$}{$12$}{$13$}{$20$}
\solbtn{1}
\taskBlock{Сюжетна задача на спільну роботу --- сесія 1.06, завдання №11}
\zadtask{У цеху на лінії виробництва пляшкової води працюють два автомати для клейового маркування. Перший маркує $14\,000$ пляшок за годину, а другий~--- $16\,000$ пляшок за годину. Обидва автомати одночасно працювали $12$ хвилин. Скільки всього пляшок було промарковано за цей час?}
\answerTable{$2000$}{$3000$}{$5000$}{$6000$}{$7500$}
\solbtn{2}
\taskBlock{Пропорційні величини (швидкість відтворення) --- сесія 3.06, завдання №3}
\zadtask{Вадим планує переглянути відеоролик тривалістю $60$ хвилин. Перед переглядом він збільшив швидкість відтворення в $1{,}5$ раза. Скільки хвилин тепер триватиме перегляд?}
\answerTable{$40$}{$90$}{$45$}{$30$}{$35$}
\solbtn{3}
\taskBlock{Дії з мішаними дробами --- сесія 3.06, завдання №7}
\zadtask{Обчисліть $5\dfrac{1}{3} : 1\dfrac{3}{5}$.}
\answerTableTall{$\dfrac{3}{10}$}{$\dfrac{128}{15}$}{$\dfrac{10}{3}$}{$\dfrac{5}{3}$}{$\dfrac{15}{128}$}
\solbtn{4}
\taskBlock{Стандартний вигляд числа --- сесія 4.06, завдання №1}
\zadtask{В астрономії для визначення відстані використовують позасистемну одиницю довжини~--- астрономічна одиниця (а.\,о.), що дорівнює середній відстані від Землі до Сонця: $1$ а.\,о. $\approx 149{,}6$~млн~\textit{км}. Запишіть це значення (у \textit{км}) у стандартному вигляді.}
\answerTable{$149{,}6 \cdot 10^6$}{$1{,}496 \cdot 10^8$}{$0{,}1496 \cdot 10^9$}{$1496 \cdot 10^5$}{$14{,}96 \cdot 10^7$}
\solbtn{5}
\taskBlock{Пропорційне обчислення (швидкість роботи) --- сесія 5.06, завдання №1}
\zadtask{Настінний принтер друкує зображення зі сталою швидкістю $3$ м$^2$ за одну годину. Скільки \textit{хвилин} принтер друкував на стіні будинку зображення площею $2{,}4$ м$^2$?}
\answerTable{$36$}{$40$}{$51$}{$48$}{$45$}
\solbtn{6}
\typeTitle{Перетворення виразів}
\taskBlock{Спрощення раціонального виразу --- сесія 23.05, завдання №7}
\zadtask{Спростіть вираз $\dfrac{(2x - 3)^2 - 9}{x}$. \nmtyear{2026}}
\answerTable{$4x - 12$}{$4x - 6$}{$4x$}{$2x - 6$}{$2x - 12$}
\solbtn{7}
\taskBlock{Спрощення раціонального виразу --- сесія 30.05, завдання №6}
\zadtask{Спростіть вираз $\dfrac{9 - x^2}{x^2 + 6x + 9}$.}
\answerTableTall{$\dfrac{3-x}{x+3}$}{$\dfrac{x-3}{x+3}$}{$\dfrac{3+x}{x-3}$}{$-\dfrac{x+3}{x-3}$}{$\dfrac{1}{x+3}$}
\solbtn{8}
\taskBlock{Спрощення виразу зі степенями --- сесія 1.06, завдання №4}
\zadtask{Спростіть вираз $\dfrac{6x^3 + 3x^3}{3}$.}
\answerTable{$3x^6$}{$6x^6$}{$5x^3$}{$3x^3$}{$7x^3$}
\solbtn{9}
\taskBlock{Розкриття дужок і спрощення виразу --- сесія 2.06, завдання №2}
\zadtask{Спростіть вираз $15 - 4(2x^2 + 3)$.}
\answerTable{$3 - 8x^2$}{$18 - 8x^2$}{$27 - 8x^2$}{$22x^2 + 33$}{$22x^2 + 3$}
\solbtn{10}
\taskBlock{Спрощення раціонального виразу --- сесія 2.06, завдання №15}
\zadtask{Спростіть вираз $\dfrac{x^2 - 4xy + 4y^2}{x - 2y} : (2y - x)$.}
\answerTable{$-1$}{$-x^2 + 4xy - y^2$}{$2y - x$}{$1$}{$x^2 - 4xy + y^2$}
\solbtn{11}
\taskBlock{Спрощення раціонального виразу --- сесія 4.06, завдання №6}
\zadtask{$\dfrac{a^2 - 25b^2}{10b + 2a} =$}
\answerTableTall{$\dfrac{a - 5b}{2}$}{$2a - 10b$}{$\dfrac{a + 5b}{2}$}{$2a + 10b$}{$\dfrac{a - 5b}{4}$}
\solbtn{12}
\taskBlock{Спрощення виразу (формули скороченого множення) --- сесія 5.06, завдання №6}
\zadtask{Спростіть вираз $(a+6)(a-6) - (a-4)^2$.}
\answerTable{$8a - 52$}{$-8a - 52$}{$8a - 20$}{$-52$}{$8a + 52$}
\solbtn{13}
\typeTitle{Тригонометрія}
\taskBlock{Спрощення тригонометричного виразу --- сесія 30.05, завдання №13}
\zadtask{Спростіть вираз $\mathrm{tg}^2 x \cdot \cos^2 x - 1$.}
\answerTableTall{$-\sin^2 x$}{$-\cos^2 x$}{$-1$}{$\cos^2 x$}{$\sin^2 x$}
\solbtn{14}
\taskBlock{Знаки тригонометричних функцій --- сесія 1.06, завдання №13}
\zadtask{Відомо, що $\cos 210^\circ \cdot \sin\alpha > 0$. Якого з наведених значень може набувати кут $\alpha$? \nmtyear{2026}}
\answerTable{$0^\circ$}{$30^\circ$}{$90^\circ$}{$150^\circ$}{$300^\circ$}
\solbtn{15}
\taskBlock{Обчислення тригонометричного виразу --- сесія 3.06, завдання №13}
\zadtask{Обчисліть значення виразу $2\sin\left(-\dfrac{\pi}{2}\right) + \cos 3\pi$.}
\answerTable{$-3$}{$-2$}{$-1$}{$2$}{$3$}
\solbtn{16}
\taskBlock{Тригонометричний вираз (формула подвійного кута) --- сесія 4.06, завдання №13}
\zadtask{Якому проміжку належить значення виразу $6 \sin 15^\circ \cos 15^\circ$?}
\answerTable{$(0;\, 1)$}{$[1;\, 2)$}{$[2;\, 3)$}{$[3;\, 5)$}{$[5;\, 8)$}
\solbtn{17}
\taskBlock{Спрощення тригонометричного виразу --- сесія 5.06, завдання №10}
\zadtask{Спростіть вираз $\mathrm{tg}^2\alpha - \mathrm{tg}^2\alpha \cdot \sin^2\alpha$.}
\answerTableTall{$\sin^2\alpha$}{$\cos^2\alpha$}{$\mathrm{tg}^2\alpha$}{$1$}{$\sin^2\alpha\cos^2\alpha$}
\solbtn{18}
\typeTitle{Логарифми}
\taskBlock{Обчислення логарифма --- сесія 23.05, завдання №9}
\zadtask{Обчисліть $\log_{0{,}5} 8$. \nmtyear{2026}}
\answerTableTall{$-2$}{$-4$}{$-3$}{$\dfrac{16}{3}$}{$3$}
\solbtn{19}
\typeTitle{Завдання на встановлення відповідності}
\taskBlock{Відповідність: вирази та їх значення --- сесія 23.05, завдання №17}
\zadtask{Установіть відповідність між виразом (1--3) та його значенням (А--Д). \nmtyear{2026}

\vspace{0.3cm}

\noindent
\begin{minipage}[t]{0.52\textwidth}
\textit{Вираз} \par \vspace{0.2cm}
\textbf{1} \ \ $\left|\dfrac{5}{3} - 2\right|$\\[0.35cm]
\textbf{2} \ \ $3^{-2} \cdot \left(\dfrac{1}{3}\right)^{-3}$\\[0.35cm]
\textbf{3} \ \ $\sqrt{3} \cdot \sin\dfrac{\pi}{3}$
\end{minipage}
\begin{minipage}[t]{0.24\textwidth}
\textit{Значення} \par \vspace{0.2cm}
\textbf{А} \ \ $\dfrac{3}{2}$\\[0.15cm]
\textbf{Б} \ \ $3$\\[0.15cm]
\textbf{В} \ \ $\dfrac{1}{3}$\\[0.15cm]
\textbf{Г} \ \ $-\dfrac{1}{3}$\\[0.15cm]
\textbf{Д} \ \ $\dfrac{1}{2}$
\end{minipage}
\hfill
\begin{minipage}[t]{0.22\textwidth}
\vspace{-0.2cm}
\begin{flushright}\matchingGrid\end{flushright}
\end{minipage}
}
\solbtn{20}
\taskBlock{Обчислення значень виразів (відповідність) --- сесія 30.05, завдання №18}
\zadtask{Установіть відповідність між виразом (1--3) і його значенням (А--Д), якщо $a = \dfrac{3}{8}$.

\vspace{0.2cm}
\noindent
\begin{minipage}[t]{0.45\textwidth}
\textit{Вираз}\par\vspace{0.15cm}
\textbf{1} \quad $\sqrt[3]{\dfrac{a}{3}}$\\[0.35cm]
\textbf{2} \quad $\log_2 a + \log_2 \dfrac{16}{3}$\\[0.35cm]
\textbf{3} \quad $\sin(4\pi a)$
\end{minipage}
\hfill
\begin{minipage}[t]{0.3\textwidth}
\textit{Значення}\par\vspace{0.15cm}
\textbf{А} \quad $2$\\[0.15cm]
\textbf{Б} \quad $1$\\[0.15cm]
\textbf{В} \quad $\dfrac{1}{2}$\\[0.2cm]
\textbf{Г} \quad $0$\\[0.15cm]
\textbf{Д} \quad $-1$
\end{minipage}
\hfill
\begin{minipage}[t]{0.2\textwidth}
\vspace{0.4cm}
\begin{flushright}\matchingGrid\end{flushright}
\end{minipage}
}
\solbtn{21}
\taskBlock{Відповідність: значення виразів (степені, логарифми) --- сесія 1.06, завдання №16}
\zadtask{Узгодьте вираз (1--3) із його значенням (А--Д), якщо $a = 4$.

\vspace{0.3cm}
\noindent
\begin{minipage}[t]{0.45\textwidth}
\textit{Вираз}\par\vspace{0.2cm}
\textbf{1} \quad $\dfrac{2a^2 - 12a}{a^2 - 36}$\\[0.5cm]
\textbf{2} \quad $0{,}04^{\,\log_5 a}$\\[0.45cm]
\textbf{3} \quad $a^0 - a^{\frac{3}{2}}$
\end{minipage}
\hfill
\begin{minipage}[t]{0.3\textwidth}
\textit{Значення}\par\vspace{0.2cm}
\textbf{А} \quad $\dfrac{1}{16}$\\[0.25cm]
\textbf{Б} \quad $-8$\\[0.2cm]
\textbf{В} \quad $-7$\\[0.2cm]
\textbf{Г} \quad $0{,}8$\\[0.2cm]
\textbf{Д} \quad $-\dfrac{1}{16}$
\end{minipage}
\hfill
\begin{minipage}[t]{0.2\textwidth}
\vspace{0.5cm}
\begin{flushright}\matchingGrid\end{flushright}
\end{minipage}
}
\solbtn{22}
\taskBlock{Відповідність: значення виразу та проміжок --- сесія 2.06, завдання №16}
\zadtask{Узгодьте вираз (1--3) із проміжком (А--Д), якому належить значення цього виразу.

\vspace{0.3cm}
\noindent
\begin{minipage}[t]{0.4\textwidth}
\textit{Вираз}\par\vspace{0.2cm}
\textbf{1} \quad $\dfrac{7}{5} - \dfrac{5}{3}$\\[0.5cm]
\textbf{2} \quad $\sqrt{3} \cdot \sqrt{5}$\\[0.4cm]
\textbf{3} \quad $\log_4 8$
\end{minipage}
\hfill
\begin{minipage}[t]{0.32\textwidth}
\textit{Проміжок}\par\vspace{0.2cm}
\textbf{А} \quad $[-1;\, 0)$\\[0.15cm]
\textbf{Б} \quad $[0;\, 1)$\\[0.15cm]
\textbf{В} \quad $[1;\, 2)$\\[0.15cm]
\textbf{Г} \quad $[2;\, 4)$\\[0.15cm]
\textbf{Д} \quad $[4;\, +\infty)$
\end{minipage}
\hfill
\begin{minipage}[t]{0.2\textwidth}
\vspace{0.5cm}
\begin{flushright}\matchingGrid\end{flushright}
\end{minipage}
}
\solbtn{23}
\taskBlock{Тотожні перетворення виразів (відповідність) --- сесія 3.06, завдання №16}
\noindent\zadnum Узгодьте вираз (1--3) з тотожно рівним йому виразом (А--Д), якщо $a \neq 0$.

\vspace{0.3cm}
\noindent
\begin{minipage}[t]{0.42\textwidth}
\textit{Вираз}\par\vspace{0.2cm}
\textbf{1} \quad $\sqrt[3]{\dfrac{27}{a^3}}$\\[0.55cm]
\textbf{2} \quad $a^{-2}\cdot a^3 \cdot \log_{27} 3$\\[0.4cm]
\textbf{3} \quad $\dfrac{(a-3)^2 - a^2 - 9}{2}$
\end{minipage}
\hfill
\begin{minipage}[t]{0.3\textwidth}
\textit{Тотожно рівний вираз}\par\vspace{0.2cm}
\textbf{А} \quad $-3a$\\[0.2cm]
\textbf{Б} \quad $\dfrac{3}{a}$\\[0.25cm]
\textbf{В} \quad $0$\\[0.2cm]
\textbf{Г} \quad $\dfrac{a}{3}$\\[0.25cm]
\textbf{Д} \quad $3a$
\end{minipage}
\hfill
\begin{minipage}[t]{0.2\textwidth}
\vspace{0.5cm}
\begin{flushright}\matchingGrid\end{flushright}
\end{minipage}
\solbtn{24}
\taskBlock{Обчислення значень виразів (відповідність) --- сесія 4.06, завдання №16}
\zadtask{Узгодьте вираз (1--3) із його значенням (А--Д), якщо $a = \dfrac{2}{3}$.

\vspace{0.3cm}
\noindent
\begin{minipage}[t]{0.4\textwidth}
\textit{Вираз}\par\vspace{0.2cm}
\textbf{1} \quad $6a$\\[0.4cm]
\textbf{2} \quad $\sqrt{\left(2^a\right)^3}$\\[0.4cm]
\textbf{3} \quad $\log_3 (1 - a)$
\end{minipage}
\hfill
\begin{minipage}[t]{0.32\textwidth}
\textit{Значення виразу}\par\vspace{0.2cm}
\textbf{А} \quad $4$\\[0.25cm]
\textbf{Б} \quad $-1$\\[0.25cm]
\textbf{В} \quad $1$\\[0.25cm]
\textbf{Г} \quad $\dfrac{20}{3}$\\[0.25cm]
\textbf{Д} \quad $2$
\end{minipage}
\hfill
\begin{minipage}[t]{0.2\textwidth}
\vspace{0.4cm}
\begin{flushright}\matchingGrid\end{flushright}
\end{minipage}
}
\solbtn{25}
\taskBlock{Обчислення значень виразів (відповідність) --- сесія 5.06, завдання №16}
\zadtask{Узгодьте вираз (1--3) із його значенням (А--Д).

\vspace{0.3cm}
\noindent
\begin{minipage}[t]{0.4\textwidth}
\textit{Вираз}\par\vspace{0.2cm}
\textbf{1} \quad $\left|\,-6 + \cos 4\pi\,\right|$\\[0.4cm]
\textbf{2} \quad $\sqrt{2}\cdot\sqrt{18}$\\[0.4cm]
\textbf{3} \quad $\log_{1/2} 32$
\end{minipage}
\hfill
\begin{minipage}[t]{0.3\textwidth}
\textit{Значення}\par\vspace{0.2cm}
\textbf{А} \quad $-6$\\[0.15cm]
\textbf{Б} \quad $-5$\\[0.15cm]
\textbf{В} \quad $5$\\[0.15cm]
\textbf{Г} \quad $6$\\[0.15cm]
\textbf{Д} \quad $7$
\end{minipage}
\hfill
\begin{minipage}[t]{0.2\textwidth}
\vspace{0.5cm}
\begin{flushright}\matchingGrid\end{flushright}
\end{minipage}}
\solbtn{26}
\typeTitle{Завдання з короткою відповіддю (відсотки)}
\taskBlock{Сюжетна задача на відсотки і частини --- сесія 23.05, завдання №21}
\zadtask{У магазині є лише музичні диски, диски з науково-популярними фільмами та диски з художніми фільмами. Кількість музичних дисків утричі більша за кількість дисків із науково-популярними фільмами й становить $40\%$ від кількості дисків з художніми фільмами. Загальна кількість дисків $345$. Визначте кількість музичних дисків у магазині. \nmtyear{2026}}
\nmtAnswerBox
\solbtn{27}
\taskBlock{Задача на відсотки (суміш) --- сесія 30.05, завдання №20}
\zadtask{Для приготування трав'яного чаю змішали м'яту й чебрець. Маса суміші становить $208$ г, причому м'яти в суміші на $60\,\%$ більше, ніж чебрецю. Скільки грамів м'яти у цій суміші?
\nmtAnswerBox}
\solbtn{28}
\taskBlock{Текстова задача на відсотки --- сесія 1.06, завдання №20}
\zadtask{У будинку $162$ квартири~--- одно-, дво- і трикімнатні. Двокімнатних квартир на $15\,\%$ менше, ніж однокімнатних, а трикімнатних~--- стільки ж, скільки двокімнатних. Скільки в будинку \textit{двокімнатних} квартир?}
\nmtAnswerBox
\solbtn{29}
\taskBlock{Сюжетна задача на відсотки (вартість таксі) --- сесія 2.06, завдання №20}
\zadtask{Перед подорожжю турист проаналізував інфографіку й дізнався вартість (у грн) поїздки від аеропорту до центру міста (або у зворотному напрямку) у деяких світових столицях. На час туристичного сезону вартість проїзду в таксі в Осло збільшується на $20\,\%$, а в Рейк'явіку~--- на $30\,\%$ від фіксованої вартості проїзду. Визначте сумарну кількість коштів, яку витратить турист на таксі від аеропорту до центру міста й назад у цих містах, якщо відвідає Осло та Рейк'явік у туристичний сезон.

\vspace{0.3cm}
\begin{center}
\begin{tikzpicture}[scale=1.0]
    % фон-рамка
    \draw[gray!40, rounded corners=3pt] (-0.3,-0.4) rectangle (9.2,5.6);
    % дані: місто / значення / довжина бара (масштаб: 5561 -> 4.2)
    \def\scale{0.000755}
    \foreach \i/\city/\val in {0/{аеропорт Лондона}/2502, 1/{аеропорт Монако}/2502, 2/{аеропорт Рейк'явіка}/2780, 3/{аеропорт Осло}/2780, 4/{аеропорт Токіо}/5561} {
        \pgfmathsetmacro\y{\i*1.05}
        \pgfmathsetmacro\barlen{\val*\scale}
        % прапорець-заглушка
        \filldraw[fill=gray!25, draw=gray!50] (0,\y+0.32) rectangle (0.45,\y+0.62);
        % назва
        \node[right, font=\small] at (0.55,\y+0.47) {\city};
        % бар
        \filldraw[fill=cyan!45, draw=cyan!60!black] (0.55,\y) rectangle (0.55+\barlen,\y+0.28);
        \node[right, font=\scriptsize\bfseries, text=white] at (0.7,\y+0.14) {\val};
    }
    % легенда
    \filldraw[fill=cyan!45, draw=cyan!60!black] (6.0,4.55) rectangle (6.35,4.85);
    \node[right, font=\scriptsize] at (6.4,4.7) {Фіксована вартість};
    \node[right, font=\scriptsize] at (6.4,4.4) {проїзду, \textit{грн}};
\end{tikzpicture}
\end{center}
\nmtAnswerBox
}
\solbtn{30}
\taskBlock{Відсотки і знижки (сюжетна задача) --- сесія 3.06, завдання №20}
\zadtask{Сім'я з трьох осіб~--- батька, мами й сина~--- запланувала сплав на байдарці. Вартість сплаву~--- $800$ грн з особи. Мамі та сину надають знижки: $20\,\%$ і $25\,\%$ відповідно (батько сплачує повну вартість). Скільки гривень заплатить сім'я за сплав?}
\nmtAnswerBox
\solbtn{31}
\taskBlock{Текстова задача на відсотки --- сесія 4.06, завдання №19}
\zadtask{Валерія за два дні прочитала $132$ вірші. За перший день вона прочитала на $20\,\%$ віршів більше, ніж за другий. Скільки віршів прочитала Валерія за перший день?
\nmtAnswerBox}
\solbtn{32}
\taskBlock{Сюжетна задача на відсотки (акційна покупка) --- сесія 5.06, завдання №19}
\zadtask{Чоловік купив три футболки за $1056$ грн. У магазині діяла акція: першу футболку продавали за повну ціну, другу~--- на $20\,\%$ дешевше, а третю~--- на $40\,\%$ дешевше за повну ціну. Скільки гривень коштувала одна футболка \textit{без} акції?}
\nmtAnswerBox
\solbtn{33}
\chapterTitle{ПОЯСНЕННЯ ДО ЗАДАЧ}
\noindent\small Стислі пояснення до всіх оригінальних задач. Натискайте \fbox{\footnotesize $\leftarrow$ Задача №$N$}, щоб повернутися.\par\vspace{0.4cm}
\par\vspace{0.3cm}\noindent\hypertarget{sol:1}{}
\begin{tcolorbox}[colback=titleBg!40, colframe=mainGreen!55, boxrule=0.4pt, arc=2pt, left=6pt, right=6pt, top=4pt, bottom=4pt, breakable]
\noindent\small\textbf{Задача №1.}\quad $0{,}5$ кг $=500$ г. Кількість порцій: $500:40=12{,}5$, тобто ціла частина $12$. Відповідь: \textbf{В} ($12$).
\par\vspace{0.1cm}\hfill\hyperlink{task:1}{\fbox{\footnotesize $\leftarrow$ Задача №1}}
\end{tcolorbox}
\par\vspace{0.3cm}\noindent\hypertarget{sol:2}{}
\begin{tcolorbox}[colback=titleBg!40, colframe=mainGreen!55, boxrule=0.4pt, arc=2pt, left=6pt, right=6pt, top=4pt, bottom=4pt, breakable]
\noindent\small\textbf{Задача №2.}\quad За годину разом маркують $14\,000+16\,000=30\,000$ пляшок. $12$ хв $=\frac{12}{60}=0{,}2$ год, тож $30\,000\cdot 0{,}2=6000$. Відповідь: \textbf{Г}.
\par\vspace{0.1cm}\hfill\hyperlink{task:2}{\fbox{\footnotesize $\leftarrow$ Задача №2}}
\end{tcolorbox}
\par\vspace{0.3cm}\noindent\hypertarget{sol:3}{}
\begin{tcolorbox}[colback=titleBg!40, colframe=mainGreen!55, boxrule=0.4pt, arc=2pt, left=6pt, right=6pt, top=4pt, bottom=4pt, breakable]
\noindent\small\textbf{Задача №3.}\quad Збільшення швидкості в $1{,}5$ раза зменшує час у стільки ж разів: $60 : 1{,}5 = 40$ (хв). Відповідь: \textbf{А} ($40$).
\par\vspace{0.1cm}\hfill\hyperlink{task:3}{\fbox{\footnotesize $\leftarrow$ Задача №3}}
\end{tcolorbox}
\par\vspace{0.3cm}\noindent\hypertarget{sol:4}{}
\begin{tcolorbox}[colback=titleBg!40, colframe=mainGreen!55, boxrule=0.4pt, arc=2pt, left=6pt, right=6pt, top=4pt, bottom=4pt, breakable]
\noindent\small\textbf{Задача №4.}\quad $5\dfrac13 : 1\dfrac35 = \dfrac{16}{3} : \dfrac{8}{5} = \dfrac{16}{3}\cdot\dfrac{5}{8} = \dfrac{10}{3}$. Відповідь: \textbf{В} ($\dfrac{10}{3}$).
\par\vspace{0.1cm}\hfill\hyperlink{task:4}{\fbox{\footnotesize $\leftarrow$ Задача №4}}
\end{tcolorbox}
\par\vspace{0.3cm}\noindent\hypertarget{sol:5}{}
\begin{tcolorbox}[colback=titleBg!40, colframe=mainGreen!55, boxrule=0.4pt, arc=2pt, left=6pt, right=6pt, top=4pt, bottom=4pt, breakable]
\noindent\small\textbf{Задача №5.}\quad $149{,}6$~млн $=149\,600\,000=1{,}496\cdot10^8$. У стандартному вигляді множник має задовольняти $1\leqslant a<10$, тож $1{,}496\cdot10^8$. Відповідь: \textbf{Б}.
\par\vspace{0.1cm}\hfill\hyperlink{task:5}{\fbox{\footnotesize $\leftarrow$ Задача №5}}
\end{tcolorbox}
\par\vspace{0.3cm}\noindent\hypertarget{sol:6}{}
\begin{tcolorbox}[colback=titleBg!40, colframe=mainGreen!55, boxrule=0.4pt, arc=2pt, left=6pt, right=6pt, top=4pt, bottom=4pt, breakable]
\noindent\small\textbf{Задача №6.}\quad Час друку: $\dfrac{2{,}4}{3}=0{,}8$ год $=0{,}8\cdot 60=48$ хв. Відповідь: \textbf{Г}.
\par\vspace{0.1cm}\hfill\hyperlink{task:6}{\fbox{\footnotesize $\leftarrow$ Задача №6}}
\end{tcolorbox}
\par\vspace{0.3cm}\noindent\hypertarget{sol:7}{}
\begin{tcolorbox}[colback=titleBg!40, colframe=mainGreen!55, boxrule=0.4pt, arc=2pt, left=6pt, right=6pt, top=4pt, bottom=4pt, breakable]
\noindent\small\textbf{Задача №7.}\quad За формулою $(a-b)^2$: $(2x-3)^2-9=4x^2-12x+9-9=4x^2-12x=4x(x-3)$. Поділивши на $x$, дістаємо $4x-12$. Відповідь: \textbf{А}.
\par\vspace{0.1cm}\hfill\hyperlink{task:7}{\fbox{\footnotesize $\leftarrow$ Задача №7}}
\end{tcolorbox}
\par\vspace{0.3cm}\noindent\hypertarget{sol:8}{}
\begin{tcolorbox}[colback=titleBg!40, colframe=mainGreen!55, boxrule=0.4pt, arc=2pt, left=6pt, right=6pt, top=4pt, bottom=4pt, breakable]
\noindent\small\textbf{Задача №8.}\quad Чисельник $9-x^2=(3-x)(3+x)$, знаменник $x^2+6x+9=(x+3)^2$. Тоді $\dfrac{(3-x)(3+x)}{(x+3)^2}=\dfrac{3-x}{x+3}$. Відповідь: \textbf{А}.
\par\vspace{0.1cm}\hfill\hyperlink{task:8}{\fbox{\footnotesize $\leftarrow$ Задача №8}}
\end{tcolorbox}
\par\vspace{0.3cm}\noindent\hypertarget{sol:9}{}
\begin{tcolorbox}[colback=titleBg!40, colframe=mainGreen!55, boxrule=0.4pt, arc=2pt, left=6pt, right=6pt, top=4pt, bottom=4pt, breakable]
\noindent\small\textbf{Задача №9.}\quad $\dfrac{6x^3+3x^3}{3}=\dfrac{9x^3}{3}=3x^3$. Відповідь: \textbf{Г}.
\par\vspace{0.1cm}\hfill\hyperlink{task:9}{\fbox{\footnotesize $\leftarrow$ Задача №9}}
\end{tcolorbox}
\par\vspace{0.3cm}\noindent\hypertarget{sol:10}{}
\begin{tcolorbox}[colback=titleBg!40, colframe=mainGreen!55, boxrule=0.4pt, arc=2pt, left=6pt, right=6pt, top=4pt, bottom=4pt, breakable]
\noindent\small\textbf{Задача №10.}\quad Розкриваємо дужки: $15-4(2x^2+3)=15-8x^2-12=3-8x^2$. Відповідь: \textbf{А}.
\par\vspace{0.1cm}\hfill\hyperlink{task:10}{\fbox{\footnotesize $\leftarrow$ Задача №10}}
\end{tcolorbox}
\par\vspace{0.3cm}\noindent\hypertarget{sol:11}{}
\begin{tcolorbox}[colback=titleBg!40, colframe=mainGreen!55, boxrule=0.4pt, arc=2pt, left=6pt, right=6pt, top=4pt, bottom=4pt, breakable]
\noindent\small\textbf{Задача №11.}\quad Чисельник $x^2-4xy+4y^2=(x-2y)^2$, тож дріб дорівнює $x-2y$. Ділення на $(2y-x)=-(x-2y)$ дає $\dfrac{x-2y}{-(x-2y)}=-1$. Відповідь: \textbf{А}.
\par\vspace{0.1cm}\hfill\hyperlink{task:11}{\fbox{\footnotesize $\leftarrow$ Задача №11}}
\end{tcolorbox}
\par\vspace{0.3cm}\noindent\hypertarget{sol:12}{}
\begin{tcolorbox}[colback=titleBg!40, colframe=mainGreen!55, boxrule=0.4pt, arc=2pt, left=6pt, right=6pt, top=4pt, bottom=4pt, breakable]
\noindent\small\textbf{Задача №12.}\quad Чисельник $a^2-25b^2=(a-5b)(a+5b)$, знаменник $10b+2a=2(a+5b)$. Скорочуємо на $(a+5b)$: $\dfrac{(a-5b)(a+5b)}{2(a+5b)}=\dfrac{a-5b}{2}$. Відповідь: \textbf{А}.
\par\vspace{0.1cm}\hfill\hyperlink{task:12}{\fbox{\footnotesize $\leftarrow$ Задача №12}}
\end{tcolorbox}
\par\vspace{0.3cm}\noindent\hypertarget{sol:13}{}
\begin{tcolorbox}[colback=titleBg!40, colframe=mainGreen!55, boxrule=0.4pt, arc=2pt, left=6pt, right=6pt, top=4pt, bottom=4pt, breakable]
\noindent\small\textbf{Задача №13.}\quad $(a+6)(a-6)-(a-4)^2 = (a^2-36)-(a^2-8a+16)=8a-52$. Відповідь: \textbf{А}.
\par\vspace{0.1cm}\hfill\hyperlink{task:13}{\fbox{\footnotesize $\leftarrow$ Задача №13}}
\end{tcolorbox}
\par\vspace{0.3cm}\noindent\hypertarget{sol:14}{}
\begin{tcolorbox}[colback=titleBg!40, colframe=mainGreen!55, boxrule=0.4pt, arc=2pt, left=6pt, right=6pt, top=4pt, bottom=4pt, breakable]
\noindent\small\textbf{Задача №14.}\quad $\mathrm{tg}^2 x\cdot\cos^2 x=\dfrac{\sin^2 x}{\cos^2 x}\cdot\cos^2 x=\sin^2 x$. Тоді $\sin^2 x-1=-(1-\sin^2 x)=-\cos^2 x$. Відповідь: \textbf{Б}.
\par\vspace{0.1cm}\hfill\hyperlink{task:14}{\fbox{\footnotesize $\leftarrow$ Задача №14}}
\end{tcolorbox}
\par\vspace{0.3cm}\noindent\hypertarget{sol:15}{}
\begin{tcolorbox}[colback=titleBg!40, colframe=mainGreen!55, boxrule=0.4pt, arc=2pt, left=6pt, right=6pt, top=4pt, bottom=4pt, breakable]
\noindent\small\textbf{Задача №15.}\quad $\cos 210^\circ<0$, тому для $\cos 210^\circ\cdot\sin\alpha>0$ потрібно $\sin\alpha<0$. Серед варіантів лише $\sin 300^\circ<0$. Відповідь: \textbf{Д}.
\par\vspace{0.1cm}\hfill\hyperlink{task:15}{\fbox{\footnotesize $\leftarrow$ Задача №15}}
\end{tcolorbox}
\par\vspace{0.3cm}\noindent\hypertarget{sol:16}{}
\begin{tcolorbox}[colback=titleBg!40, colframe=mainGreen!55, boxrule=0.4pt, arc=2pt, left=6pt, right=6pt, top=4pt, bottom=4pt, breakable]
\noindent\small\textbf{Задача №16.}\quad $\sin\left(-\dfrac{\pi}{2}\right) = -1$, $\cos 3\pi = -1$, тому $2\cdot(-1) + (-1) = -3$. Відповідь: \textbf{А} ($-3$).
\par\vspace{0.1cm}\hfill\hyperlink{task:16}{\fbox{\footnotesize $\leftarrow$ Задача №16}}
\end{tcolorbox}
\par\vspace{0.3cm}\noindent\hypertarget{sol:17}{}
\begin{tcolorbox}[colback=titleBg!40, colframe=mainGreen!55, boxrule=0.4pt, arc=2pt, left=6pt, right=6pt, top=4pt, bottom=4pt, breakable]
\noindent\small\textbf{Задача №17.}\quad За формулою $2\sin\alpha\cos\alpha=\sin2\alpha$: $6\sin15^\circ\cos15^\circ=3\sin30^\circ=3\cdot\dfrac{1}{2}=1{,}5\in[1;\,2)$. Відповідь: \textbf{Б}.
\par\vspace{0.1cm}\hfill\hyperlink{task:17}{\fbox{\footnotesize $\leftarrow$ Задача №17}}
\end{tcolorbox}
\par\vspace{0.3cm}\noindent\hypertarget{sol:18}{}
\begin{tcolorbox}[colback=titleBg!40, colframe=mainGreen!55, boxrule=0.4pt, arc=2pt, left=6pt, right=6pt, top=4pt, bottom=4pt, breakable]
\noindent\small\textbf{Задача №18.}\quad $\mathrm{tg}^2\alpha-\mathrm{tg}^2\alpha\sin^2\alpha=\mathrm{tg}^2\alpha(1-\sin^2\alpha)=\dfrac{\sin^2\alpha}{\cos^2\alpha}\cdot\cos^2\alpha=\sin^2\alpha$. Відповідь: \textbf{А}.
\par\vspace{0.1cm}\hfill\hyperlink{task:18}{\fbox{\footnotesize $\leftarrow$ Задача №18}}
\end{tcolorbox}
\par\vspace{0.3cm}\noindent\hypertarget{sol:19}{}
\begin{tcolorbox}[colback=titleBg!40, colframe=mainGreen!55, boxrule=0.4pt, arc=2pt, left=6pt, right=6pt, top=4pt, bottom=4pt, breakable]
\noindent\small\textbf{Задача №19.}\quad Оскільки $0{,}5=2^{-1}$, а $8=2^3$, то $\log_{0{,}5}8=\dfrac{\log_2 8}{\log_2 0{,}5}=\dfrac{3}{-1}=-3$. Відповідь: \textbf{В}.
\par\vspace{0.1cm}\hfill\hyperlink{task:19}{\fbox{\footnotesize $\leftarrow$ Задача №19}}
\end{tcolorbox}
\par\vspace{0.3cm}\noindent\hypertarget{sol:20}{}
\begin{tcolorbox}[colback=titleBg!40, colframe=mainGreen!55, boxrule=0.4pt, arc=2pt, left=6pt, right=6pt, top=4pt, bottom=4pt, breakable]
\noindent\small\textbf{Задача №20.}\quad $1)\ \left|\dfrac{5}{3}-2\right|=\left|-\dfrac{1}{3}\right|=\dfrac{1}{3}$~(В); $2)\ 3^{-2}\cdot 3^{3}=3^{1}=3$~(Б); $3)\ \sqrt{3}\cdot\sin\dfrac{\pi}{3}=\sqrt{3}\cdot\dfrac{\sqrt{3}}{2}=\dfrac{3}{2}$~(А). Відповідь: \textbf{1--В, 2--Б, 3--А}.
\par\vspace{0.1cm}\hfill\hyperlink{task:20}{\fbox{\footnotesize $\leftarrow$ Задача №20}}
\end{tcolorbox}
\par\vspace{0.3cm}\noindent\hypertarget{sol:21}{}
\begin{tcolorbox}[colback=titleBg!40, colframe=mainGreen!55, boxrule=0.4pt, arc=2pt, left=6pt, right=6pt, top=4pt, bottom=4pt, breakable]
\noindent\small\textbf{Задача №21.}\quad $a=\dfrac38$. \textbf{1}: $\sqrt[3]{\dfrac{a}{3}}=\sqrt[3]{\dfrac18}=\dfrac12$ (В). \textbf{2}: $\log_2\!\left(a\cdot\dfrac{16}{3}\right)=\log_2 2=1$ (Б). \textbf{3}: $\sin\!\left(4\pi\cdot\dfrac38\right)=\sin\dfrac{3\pi}{2}=-1$ (Д). Відповідь: \textbf{1~--~В, 2~--~Б, 3~--~Д}.
\par\vspace{0.1cm}\hfill\hyperlink{task:21}{\fbox{\footnotesize $\leftarrow$ Задача №21}}
\end{tcolorbox}
\par\vspace{0.3cm}\noindent\hypertarget{sol:22}{}
\begin{tcolorbox}[colback=titleBg!40, colframe=mainGreen!55, boxrule=0.4pt, arc=2pt, left=6pt, right=6pt, top=4pt, bottom=4pt, breakable]
\noindent\small\textbf{Задача №22.}\quad При $a=4$: \textbf{1)} $\dfrac{2a^2-12a}{a^2-36}=\dfrac{2a(a-6)}{(a-6)(a+6)}=\dfrac{2a}{a+6}=\dfrac{8}{10}=0{,}8$ (Г); \textbf{2)} $0{,}04^{\log_5 4}=(5^{-2})^{\log_5 4}=4^{-2}=\dfrac{1}{16}$ (А); \textbf{3)} $a^0-a^{3/2}=1-8=-7$ (В). Відповідь: \textbf{1\,--\,Г, 2\,--\,А, 3\,--\,В}.
\par\vspace{0.1cm}\hfill\hyperlink{task:22}{\fbox{\footnotesize $\leftarrow$ Задача №22}}
\end{tcolorbox}
\par\vspace{0.3cm}\noindent\hypertarget{sol:23}{}
\begin{tcolorbox}[colback=titleBg!40, colframe=mainGreen!55, boxrule=0.4pt, arc=2pt, left=6pt, right=6pt, top=4pt, bottom=4pt, breakable]
\noindent\small\textbf{Задача №23.}\quad $\dfrac{7}{5}-\dfrac{5}{3}=\dfrac{21-25}{15}=-\dfrac{4}{15}\in[-1;0)$ -- \textbf{А}; $\sqrt{3}\cdot\sqrt{5}=\sqrt{15}\approx3{,}87\in[2;4)$ -- \textbf{Г}; $\log_4 8=\dfrac{3}{2}=1{,}5\in[1;2)$ -- \textbf{В}. Відповідь: \textbf{1--А, 2--Г, 3--В}.
\par\vspace{0.1cm}\hfill\hyperlink{task:23}{\fbox{\footnotesize $\leftarrow$ Задача №23}}
\end{tcolorbox}
\par\vspace{0.3cm}\noindent\hypertarget{sol:24}{}
\begin{tcolorbox}[colback=titleBg!40, colframe=mainGreen!55, boxrule=0.4pt, arc=2pt, left=6pt, right=6pt, top=4pt, bottom=4pt, breakable]
\noindent\small\textbf{Задача №24.}\quad \textbf{1}: $\sqrt[3]{\dfrac{27}{a^3}} = \dfrac{3}{a}$~--- Б. \textbf{2}: $a^{-2}\cdot a^3\cdot\log_{27}3 = a\cdot\dfrac13 = \dfrac{a}{3}$~--- Г. \textbf{3}: $\dfrac{(a-3)^2 - a^2 - 9}{2} = \dfrac{-6a}{2} = -3a$~--- А. Відповідь: \textbf{1~--~Б, 2~--~Г, 3~--~А}.
\par\vspace{0.1cm}\hfill\hyperlink{task:24}{\fbox{\footnotesize $\leftarrow$ Задача №24}}
\end{tcolorbox}
\par\vspace{0.3cm}\noindent\hypertarget{sol:25}{}
\begin{tcolorbox}[colback=titleBg!40, colframe=mainGreen!55, boxrule=0.4pt, arc=2pt, left=6pt, right=6pt, top=4pt, bottom=4pt, breakable]
\noindent\small\textbf{Задача №25.}\quad При $a=\dfrac{2}{3}$: $6a=4$ (А); $\sqrt{(2^a)^3}=\sqrt{2^{3a}}=\sqrt{2^2}=2$ (Д); $\log_3(1-a)=\log_3\dfrac{1}{3}=-1$ (Б). Відповідь: \textbf{1~--~А, 2~--~Д, 3~--~Б}.
\par\vspace{0.1cm}\hfill\hyperlink{task:25}{\fbox{\footnotesize $\leftarrow$ Задача №25}}
\end{tcolorbox}
\par\vspace{0.3cm}\noindent\hypertarget{sol:26}{}
\begin{tcolorbox}[colback=titleBg!40, colframe=mainGreen!55, boxrule=0.4pt, arc=2pt, left=6pt, right=6pt, top=4pt, bottom=4pt, breakable]
\noindent\small\textbf{Задача №26.}\quad $1$: $|-6+\cos4\pi|=|-6+1|=5$ (В); $2$: $\sqrt{2}\cdot\sqrt{18}=\sqrt{36}=6$ (Г); $3$: $\log_{1/2}32=\log_{1/2}2^5=-5$ (Б). Відповідь: \textbf{1~--~В, 2~--~Г, 3~--~Б}.
\par\vspace{0.1cm}\hfill\hyperlink{task:26}{\fbox{\footnotesize $\leftarrow$ Задача №26}}
\end{tcolorbox}
\par\vspace{0.3cm}\noindent\hypertarget{sol:27}{}
\begin{tcolorbox}[colback=titleBg!40, colframe=mainGreen!55, boxrule=0.4pt, arc=2pt, left=6pt, right=6pt, top=4pt, bottom=4pt, breakable]
\noindent\small\textbf{Задача №27.}\quad Нехай науково-популярних~--- $x$, тоді музичних~--- $3x$, а художніх~--- $\dfrac{3x}{0{,}4}=7{,}5x$. Сума: $x+3x+7{,}5x=11{,}5x=345$, звідки $x=30$. Музичних: $3x=90$. Відповідь: \textbf{90}.
\par\vspace{0.1cm}\hfill\hyperlink{task:27}{\fbox{\footnotesize $\leftarrow$ Задача №27}}
\end{tcolorbox}
\par\vspace{0.3cm}\noindent\hypertarget{sol:28}{}
\begin{tcolorbox}[colback=titleBg!40, colframe=mainGreen!55, boxrule=0.4pt, arc=2pt, left=6pt, right=6pt, top=4pt, bottom=4pt, breakable]
\noindent\small\textbf{Задача №28.}\quad Нехай чебрецю $c$ г, тоді м'яти $1{,}6c$ г. Маємо $c+1{,}6c=208$, звідки $2{,}6c=208$, $c=80$. М'яти: $1{,}6\cdot80=128$ г. Відповідь: $128$.
\par\vspace{0.1cm}\hfill\hyperlink{task:28}{\fbox{\footnotesize $\leftarrow$ Задача №28}}
\end{tcolorbox}
\par\vspace{0.3cm}\noindent\hypertarget{sol:29}{}
\begin{tcolorbox}[colback=titleBg!40, colframe=mainGreen!55, boxrule=0.4pt, arc=2pt, left=6pt, right=6pt, top=4pt, bottom=4pt, breakable]
\noindent\small\textbf{Задача №29.}\quad Нехай однокімнатних $x$, тоді дво- і трикімнатних по $0{,}85x$. Усього $x+0{,}85x+0{,}85x=2{,}7x=162$, звідки $x=60$. Двокімнатних: $0{,}85\cdot 60=51$. Відповідь: \textbf{$51$}.
\par\vspace{0.1cm}\hfill\hyperlink{task:29}{\fbox{\footnotesize $\leftarrow$ Задача №29}}
\end{tcolorbox}
\par\vspace{0.3cm}\noindent\hypertarget{sol:30}{}
\begin{tcolorbox}[colback=titleBg!40, colframe=mainGreen!55, boxrule=0.4pt, arc=2pt, left=6pt, right=6pt, top=4pt, bottom=4pt, breakable]
\noindent\small\textbf{Задача №30.}\quad Фіксована вартість в Осло й Рейк'явіку по $2780$ грн. У сезон: Осло $2780\cdot1{,}2$, Рейк'явік $2780\cdot1{,}3$, а туди й назад -- подвоюємо: $2\cdot(2780\cdot1{,}2+2780\cdot1{,}3)=2\cdot2780\cdot2{,}5=13900$ грн.
\par\vspace{0.1cm}\hfill\hyperlink{task:30}{\fbox{\footnotesize $\leftarrow$ Задача №30}}
\end{tcolorbox}
\par\vspace{0.3cm}\noindent\hypertarget{sol:31}{}
\begin{tcolorbox}[colback=titleBg!40, colframe=mainGreen!55, boxrule=0.4pt, arc=2pt, left=6pt, right=6pt, top=4pt, bottom=4pt, breakable]
\noindent\small\textbf{Задача №31.}\quad Батько платить повну вартість $800$ грн, мама $800\cdot0{,}8 = 640$ грн, син $800\cdot0{,}75 = 600$ грн. Усього $800 + 640 + 600 = 2040$ грн. Відповідь: \textbf{$2040$}.
\par\vspace{0.1cm}\hfill\hyperlink{task:31}{\fbox{\footnotesize $\leftarrow$ Задача №31}}
\end{tcolorbox}
\par\vspace{0.3cm}\noindent\hypertarget{sol:32}{}
\begin{tcolorbox}[colback=titleBg!40, colframe=mainGreen!55, boxrule=0.4pt, arc=2pt, left=6pt, right=6pt, top=4pt, bottom=4pt, breakable]
\noindent\small\textbf{Задача №32.}\quad Нехай за другий день прочитано $x$ віршів, тоді за перший $1{,}2x$. Сума: $x+1{,}2x=132$, $2{,}2x=132$, $x=60$. За перший день $1{,}2\cdot60=72$. Відповідь: \textbf{72}.
\par\vspace{0.1cm}\hfill\hyperlink{task:32}{\fbox{\footnotesize $\leftarrow$ Задача №32}}
\end{tcolorbox}
\par\vspace{0.3cm}\noindent\hypertarget{sol:33}{}
\begin{tcolorbox}[colback=titleBg!40, colframe=mainGreen!55, boxrule=0.4pt, arc=2pt, left=6pt, right=6pt, top=4pt, bottom=4pt, breakable]
\noindent\small\textbf{Задача №33.}\quad Нехай повна ціна футболки $x$ грн. Тоді $x+0{,}8x+0{,}6x=2{,}4x=1056$, звідки $x=440$. Відповідь: \textbf{$440$}.
\par\vspace{0.1cm}\hfill\hyperlink{task:33}{\fbox{\footnotesize $\leftarrow$ Задача №33}}
\end{tcolorbox}
\chapterTitle{ВІДПОВІДІ}
\begin{multicols}{3}\noindent
\textbf{1.}~В\\
\textbf{2.}~Г\\
\textbf{3.}~А\\
\textbf{4.}~В\\
\textbf{5.}~Б\\
\textbf{6.}~Г\\
\textbf{7.}~А\\
\textbf{8.}~А\\
\textbf{9.}~Г\\
\textbf{10.}~А\\
\textbf{11.}~А\\
\textbf{12.}~А\\
\textbf{13.}~А\\
\textbf{14.}~Б\\
\textbf{15.}~Д\\
\textbf{16.}~А\\
\textbf{17.}~Б\\
\textbf{18.}~А\\
\textbf{19.}~В\\
\textbf{20.}~1~--~В, 2~--~Б, 3~--~А\\
\textbf{21.}~1~--~В, 2~--~Б, 3~--~Д\\
\textbf{22.}~1~--~Г, 2~--~А, 3~--~В\\
\textbf{23.}~1~--~А, 2~--~Г, 3~--~В\\
\textbf{24.}~1~--~Б, 2~--~Г, 3~--~А\\
\textbf{25.}~1~--~А, 2~--~Д, 3~--~Б\\
\textbf{26.}~1~--~В, 2~--~Г, 3~--~Б\\
\textbf{27.}~$90$\\
\textbf{28.}~$128$\\
\textbf{29.}~$51$\\
\textbf{30.}~$13900$\\
\textbf{31.}~$2040$\\
\textbf{32.}~$72$\\
\textbf{33.}~$440$\\
\end{multicols}
\end{document}
